% ==== AUTO-GENERADA — no editar a mano (regenerar con generate_thesis_tables.py) ====
\begin{table}[htbp]
\centering
\caption[Interpolación a un mismo tamaño (heavy-hex)]{Interpolación a un mismo tamaño en heavy-hex: la GNN reproduce $\theta^*(h)$ en valores de $h$ no vistos, para cada $N$ de entrenamiento.}
\label{tab:auto_heavy_hex_intra_n}
\begin{tabular}{rrrr}
\toprule
$N$ & $|\Delta E|$ & gap & $\Delta E/\mathrm{gap}$ (\%) \\
\midrule
8 & 0,1865 & 5,4983 & 3,11 \\
10 & 0,0870 & 5,4271 & 1,47 \\
14 & 0,0609 & 4,8567 & 1,83 \\
16 & 0,0860 & 5,3333 & 1,39 \\
18 & 1,2439 & 5,5377 & 31,96 \\
20 & 0,0211 & 4,2064 & 0,69 \\
\bottomrule
\end{tabular}
\\[2pt]
\footnotesize Evaluado sobre la grilla de $h$ canónica del régimen de extrapolación, $h \in \{2,5; 3,0; 3,5; 4,0; 4,5; 5,0\}$, restringida al régimen válido ($h \geq h_{\min}$) de cada $N$; por eso los $N$ mayores contribuyen con menos puntos. $|\Delta E|$ y gap son medias sobre esos puntos, en unidades de $J = 1$. La columna $\Delta E/\mathrm{gap}$ es la media de los cocientes \emph{por punto} (no el cociente de las dos columnas anteriores): cerca de la frontera algún punto tiene el gap estrecho y su cociente individual es grande, de modo que la media de cocientes puede superar al cociente de las medias. La métrica primaria es $|\Delta E|$; $\Delta E/\mathrm{gap}$ es el criterio de clasificación normalizado.
\end{table}
